\documentclass[11pt,a4paper]{article}
\usepackage[margin=27mm]{geometry}
\usepackage{amsmath,amssymb,amsthm,fontspec,unicode-math,hyperref,booktabs,longtable}
\setmainfont{STIXGeneral}
\setmathfont{STIX Two Math}
\setmonofont{Menlo}
\usepackage{array,calc,fvextra}
\usepackage{newunicodechar}
\newunicodechar{∎}{\ensuremath{\blacksquare}}
\hypersetup{hidelinks}
\setlength{\emergencystretch}{2em}
\providecommand{\tightlist}{\setlength{\itemsep}{0pt}\setlength{\parskip}{0pt}}
\newcommand{\gen}[1]{\langle #1\rangle}
\newcommand{\Sym}{\operatorname{Sym}}
\newcommand{\Alt}{\operatorname{Alt}}
\newcommand{\core}{\operatorname{core}}
\newcommand{\F}{\mathbb F}
\newtheorem{theorem}{Theorem}
\newtheorem{lemma}[theorem]{Lemma}
\newtheorem{proposition}[theorem]{Proposition}
\newtheorem{corollary}[theorem]{Corollary}

\newcommand{\Aut}{\operatorname{Aut}}
\newcommand{\Fix}{\operatorname{Fix}}
\newcommand{\Y}{\mathsf Y}

\title{Forest coordinates and relative generation}
\author{Human Compression project}\date{Unified five-application draft}
\begin{document}\maketitle
\begin{abstract}
We develop common generation, subgroup and flag arguments for five
families of incidence geometries. A laminar extension theorem preserves
a prescribed local map, so ordinary tree flags, parity conditions and
four-point decorations use the same complete extension proof.
Normal-quotient normalization and binary-module bounds give a single
extremal symmetric-group classification used by two applications.
We also retain complete primitive and imprimitive bounds, finite-core
intersections and exact tree presentations needed by the high-rank
converse. The local Frattini and graph-automorphism arguments remain
in their applications. All source hypotheses and retained outside
results are explicit; this edition makes no claim about its measured
length before independent review.
\end{abstract}
\section{Introduction}
There are two different directions in these results. Once a tree or a
finite decorated core is known, subgroup intersections and completion
of flags can be proved by common constructions. Proving that every
allowed marking has that form is a separate classification problem.
We keep those directions separate and give complete proofs of both
where the applications require them.

The common coset and relative-generation interfaces serve the original
four applications. The laminar extension theorem strengthens their
forest proof by retaining a prescribed map on a root set. The three
four-point decorations then require only their actual local verification.
The ordinary/sign extremal classification is proved once here, with
Cameron--Cara and Whiston attribution, and supplies the involution-only
maximum-rank classification as a corollary. The all-degree transitive
parabolic exclusion remains a substantial additional argument; it is
currently used by the high-rank application alone and is included in full.

A C-group marking is a finite generating family of distinct nonidentity
involutions satisfying the intersection property for all subsets.
No string commutations are assumed. Independence bounds below concern
arbitrary group elements and do not require involutions or the
intersection property. Finite-group and finite-degree assumptions
are stated where needed.
\section{Flags, parabolics and coset coordinates}
All type sets are finite. Groups need not be finite unless this is stated.
A typed incidence system has a reflexive symmetric incidence relation,
with distinct elements of the same type never incident. A flag is a
pairwise incident set. A chamber has every type. A geometry is such a
system in which every flag extends to a chamber. It is firm if every
nonchamber flag has at least two chamber extensions. A residue consists of
the elements incident with a fixed flag and having a type absent from
that flag. Residual connectedness means that every residue of rank at
least two has a connected incidence graph. The specified group action
is part of the data.

For subgroups $H_i\leq G$, indexed by $I$, write
\[
 H_J=\bigcap_{j\in J}H_j,\qquad H_\varnothing=G,\qquad B=H_I.
\]
The elements of type $i$ are the left cosets of $H_i$, and two cosets of
different types are incident if they intersect. The action is left
multiplication. Inversion converts this convention to the right cosets
of the sources. We will use the following three conditions:
\begin{description}
\item[(G1)] The indexed intersections $H_J$ are all distinct.
\item[(G2)] If $|I\setminus J|\geq2$, then
$H_J=\gen{H_{J\cup\{k\}}:k\in I\setminus J}$.
\item[(G3)] Every family of cosets of different types with pairwise
nonempty intersections has a common element.
\end{description}

\noindent\textbf{Proposition UB.Coset (the complete coset interface).}
Conditions (G1)--(G3) hold if and only if the coset system is firm,
residually connected and flag-transitive under its given action.
The residue at the standard flag of type $J$ is the coset geometry of
$H_J$ with subgroups $H_{J\cup\{k\}}$, $k\notin J$.

\begin{proof}
Under (G3), a flag has a representative $g$ common to its cosets and is
therefore contained in the chamber $(gH_i:i\in I)$. Translating $g$ to
the identity shows flag transitivity. Elements of type $k$ incident with
the standard flag of type $J$ can, again by (G3), be represented by an
element of $H_J$. Two such representatives give the same coset precisely
when they differ by $H_J\cap H_k$. This proves the residue description.
Condition (G1) gives $H_{I\setminus\{i\}}>B$, so every panel has at least
two elements. Every shorter flag is contained in a panel; this proves
firmness.

For completeness, a coset incidence graph for subgroups $(L_i:i\in S)$
of a group $L$, with $|S|\geq2$, is connected exactly when the $L_i$
generate $L$. Every edge can be represented by a common group element.
Successive representatives along a path differ by elements of the
subgroup corresponding to their shared vertex. Thus a connected path
from the identity cosets to $gL_i$ places $g$ in $\gen{L_j:j\in S}$.
Conversely, for $a\in L_j$ and $i\ne j$, the vertices
$L_i,L_j=aL_j,aL_i$ give a path. Translating and concatenating these
paths for a word in the $L_j$ connects every coset to the standard ones;
the standard vertices are mutually incident. Applying this criterion
to the described residues proves residual connectedness from (G2).

Conversely, flag transitivity takes any pairwise incident family to the
standard flag of its types, so it has a common representative and (G3)
holds. Firmness and panel transitivity give an element
$s_i\in H_{I\setminus\{i\}}\setminus B$ for every $i$. Its intersection
with $H_i$ is $B$, hence $s_i\notin H_i$. If $J\ne K$, choosing a type
in one set difference shows that $s_i$ belongs to one of $H_J,H_K$ and
not to the other. This proves (G1). Residual connectedness, the residue
description and the graph criterion just proved give (G2).
\end{proof}

An action is primitive when it is transitive and has no nontrivial
invariant partition. A geometry with its group action is PRI if the
action on every type is primitive, and WPRI if the action on at least
one type is primitive. The residual versions, RPRI and RWPRI, apply
to each residue of positive rank, including the whole geometry. In coset
coordinates, RWPRI says
\[
\tag{G4}
J\subsetneq I\quad\Longrightarrow\quad
\text{some }H_{J\cup\{k\}}\text{ is maximal in }H_J,
\quad k\notin J.
\]
This is the elementary equivalence between primitivity on a coset space
and maximality of its stabilizer. It refers to the stabilizer's induced
action; a kernel of that action does not change its invariant partitions.

\noindent\textbf{Lemma UB.Flags (residual weak primitivity gives flags).}
A finite-rank geometry with a residually weakly primitive group action
is flag-transitive.
\begin{proof}
Fix a chamber. In the empty-flag residue choose a type on which the group
is primitive, hence transitive. Match the element of that type in any
second chamber to the fixed one. In the residue of the matched flag,
choose a remaining primitive type and use its flag stabilizer to match
the next element. At most $|I|$ steps match the chambers. Since every
flag extends to a chamber, transitivity follows for flags of every
fixed type. The typewise bijection from cosets of stabilizers to
elements preserves incidence: transitivity on pairs of incident
elements supplies a common representative.
\end{proof}

For $B\leq G$, a finite indexed family $(s_i:i\in I)$ is independent
relative to $B$ if
$s_i\notin\gen{B,s_j:j\ne i}$ for every $i$. If it also generates
$G$ together with $B$, it is a relative generating family. For finite
$G$ let $\mu(G,B)$ and $\mu'(G,B)$ be the largest cardinalities of
relative independent generating families and relative independent
families, respectively. Omit $B$ if $B=1$. Independence for a group
acting on another group means independence under generation by all
images of the other family members; it is denoted $\mu'_G(N)$ in the
normal-quotient application.

\section{Relative generation from residual maximality}
For a finite indexed subgroup family \(H_i\leq G\), write \(H_J=\bigcap_{j\in J}H_j\), with \(H_\varnothing=G\) and \(B=H_I\).

\textbf{Theorem UB.Relative (generation from residual maximality).} Suppose all \(H_J\), \(J\subseteq I\), are distinct, and for every \(J\subsetneq I\) there is \(k\notin J\) such that \(H_{J\cup\{k\}}\) is maximal in \(H_J\). Choose \begin{equation}
s_i\in H_{I\setminus\{i\}}\setminus H_i.
\label{eq:02-relative:6-2}
\end{equation} Then \begin{equation}
H_J=\langle B,s_i:i\notin J\rangle\quad(J\subseteq I).
\label{eq:02-relative:6-3}
\end{equation} In particular these elements generate \(G\) irredundantly relative to \(B\), and their relative parabolics satisfy \begin{equation}
\langle B,s_i:i\in S\rangle\cap\langle B,s_i:i\in T\rangle
=\langle B,s_i:i\in S\cap T\rangle.
\label{eq:02-relative:6-4}
\end{equation} The group need not be finite.

\emph{Proof.} Distinctness of \(H_{I\setminus\{i\}}\) and \(B\) makes \eqref{eq:02-relative:6-2} possible, since its intersection with \(H_i\) is \(B\). Every \(s_j\) with \(j\ne i\), and every element of \(B\), belongs to \(H_i\), while \(s_i\) does not. This proves relative independence.

Prove \eqref{eq:02-relative:6-3} by downward induction from \(J=I\). Choose \(k\) as in the hypothesis. By induction, \(H_{J\cup\{k\}}\) is generated by \(B\) and the \(s_i\) with \(i\notin J\cup\{k\}\). The element \(s_k\) belongs to \(H_J\), because \(J\subseteq I\setminus\{k\}\), and lies outside \(H_{J\cup\{k\}}\), because it lies outside \(H_k\). Maximality therefore says that adjoining \(s_k\) generates \(H_J\). This is \eqref{eq:02-relative:6-3}. Finally substitute \(J=I\setminus S\) and \(I\setminus T\) into the elementary equality \(H_J\cap H_{J'}=H_{J\cup J'}\) to obtain \eqref{eq:02-relative:6-4}. \(\square\)

\noindent\textbf{Theorem UB.Generation (the geometric interface).}
For a finite-rank firm geometry with a residually weakly primitive action,
choose one chamber and its type stabilizers $H_i$. Every choice
\[
 s_i\in H_{I\setminus\{i\}}\setminus H_i
\]
is independent relative to $B=H_I$, generates $G$ over $B$, and satisfies
\[
 H_J=\gen{B,s_i:i\notin J}\qquad(J\subseteq I).
\]
Every such geometry is residually connected. For finite $G$ its rank is
at most $\mu(G,B)$. Without residual weak primitivity, a firm
flag-transitive geometry still has a relatively independent family of
panel choices and, for finite $G$, rank at most $\mu'(G,B)$.
\begin{proof}
Lemma UB.Flags supplies flag transitivity and coset coordinates for the
first assertion. Firmness gives distinct indexed intersections by the
panel argument in Proposition UB.Coset. Residual weak primitivity is
exactly (G4), so Theorem UB.Relative applies to the displayed choices.
Its formula implies (G2): if at least two types remain, each of the
remaining generators belongs to an intersection obtained by fixing
another remaining type. Consequently Proposition UB.Coset gives residual
connectedness. The relative generating family has size equal to the
rank, proving the bound by $\mu(G,B)$.

For the last assertion no maximality is needed. Firmness and flag
transitivity again give panel choices. For a fixed $i$, all other choices
and $B$ lie in $H_i$, whereas $s_i$ does not. They are therefore relatively
independent, which gives the bound by $\mu'(G,B)$.
\end{proof}

\section{Extending prescribed maps through laminar flags}
A family of subsets is laminar if any two are nested or disjoint.
The following statement retains an arbitrary prescribed local map;
that is the part not stated in the original UB.F interface.

\begin{theorem}[UB.Laminar]\label{thm:laminar}
Let $\Omega$ be finite, let $Q\subseteq\Omega$, and let
$(A_i)_{i\in I}$ be a finite indexed laminar family of subsets of
$\Omega\setminus Q$. Choose $J\subseteq I$ and subsets $X_i$,
$i\in J$, with
\[
 |X_i|=|A_i|,\qquad |X_i\cap X_j|=|A_i\cap A_j|.
\]
Write $D=\Omega\setminus\bigcup_{i\in J}X_i$.
Every injection $u:Q\to D$ extends to a permutation $g$ of $\Omega$
satisfying $gA_i=X_i$ for all $i\in J$.
\end{theorem}
\begin{proof}
If a selected set is $\Omega$, then $Q=\varnothing$ and its target
is necessarily $\Omega$. Discard that vacuous constraint; the prescribed
map is empty in this case. We may therefore assume that every selected
set is proper, using the complement of the remaining targets as the
root region.
The intersection cardinalities enforce the same disjointness and
containment relations among the $X_i$ as among the $A_i$. Equal
standard sets give equal target sets and can be treated once. Adjoin
$\Omega$ above all the selected sets. For each selected container,
subtract its maximal selected proper subcontainers. Since those
subcontainers are disjoint, the difference has size equal to the
container's size minus the sum of their sizes. The corresponding
standard and target differences therefore have equal cardinalities.
They partition $\Omega$ on each side. The root difference contains
$Q$ on the standard side and is exactly $D$ on the target side.
Extend $u$ to a bijection of these two differences, and choose arbitrary
bijections on every other corresponding pair of differences. Their
union is the required permutation.
\end{proof}


\begin{lemma}[UB.Character]\label{lem:character}
Suppose at least one permutation realizes prescribed laminar images and
a prescribed partial bijection. Fix one such permutation $g_0$, and
let $K$ be the subgroup preserving every standard selected subset and
fixing every point in the domain of that partial bijection.
All realizations are exactly $g_0K$. For a homomorphism
$\chi:S_\Omega\to C$ and $c\in C$, a realization with $\chi(g)=c$
exists exactly when $\chi(g_0)^{-1}c\in\chi(K)$.
\end{lemma}
\begin{proof}
Two realizations differ on the right by a permutation satisfying exactly
the defining conditions of $K$, and every such right multiplication
preserves the prescribed data. Applying $\chi$ proves the second claim.
\end{proof}


These statements make explicit the extension freedom in the familiar atom-matching proof. Their use below does not claim a new historical classification or a prior-art determination.

\section{Weighted forests and complete flags}
Let \(T\) be a finite tree with vertex set \(V\) and edge set \(E\). Give each vertex a finite nonempty block \(B_v\), with the blocks pairwise disjoint, and put \(\Omega=\bigsqcup_vB_v\). For \(F\subseteq E\), let \begin{equation}
K_F=\prod_C S_{\bigcup_{v\in C}B_v},
\label{eq:03-forest:2-1}
\end{equation} where \(C\) runs over the components of the forest \((V,F)\). Root \(T\), and let \(A_e\subseteq\Omega\) be the union of blocks on the side of \(T-e\) away from the root. In the finite symmetric group \(S_\Omega\), set \(H_e=K_{E\setminus\{e\}}=S_{A_e}\times S_{\Omega\setminus A_e}\).

\textbf{Theorem UB.F (forest intersections and flag completion).} The following hold.

\begin{enumerate}
\def\labelenumi{\arabic{enumi}.}
\tightlist
\item
  \(K_F\cap K_D=K_{F\cap D}\). Moreover \(K_F\) is generated by \(K_\varnothing\) together with the groups \(K_{\{e\}}\), \(e\in F\).
\item
  The coset system of the \(H_e\) is a firm, residually connected, flag-transitive geometry. Its type-\(e\) elements identify with the \(|A_e|\)-subsets of \(\Omega\), and distinct types satisfy \begin{equation}
  X_e\mathrel{*}X_f\quad\Longleftrightarrow\quad
  |X_e\cap X_f|=|A_e\cap A_f|.
  \label{eq:03-forest:2-2}
  \end{equation}
\item
  Its Borel subgroup is \(K_\varnothing=\prod_vS_{B_v}\). The residue whose remaining types are \(F\) is the join of the same constructions on the nontrivial components of \((V,F)\). Its acting group is \(K_F\); components with no edges contribute only a kernel on this residue.
\end{enumerate}

\emph{Proof.} A permutation belongs to both groups in (1) exactly when it preserves every intersection of their component blocks. Two vertices lie in a common block of both forests precisely when their unique tree path belongs to both edge sets. Thus the common refinement is the component partition for \(F\cap D\). Transpositions inside the blocks, together with transpositions joining adjacent blocks along retained edges, connect each component; they generate its full symmetric group. This proves (1).

The coset \(gH_e\) maps to \(gA_e\). Two pairs of subsets are conjugate under \(S_\Omega\) exactly when their four Boolean atoms have the same sizes, giving \eqref{eq:03-forest:2-2}. The sets \(A_e\) are laminar: any two are nested or disjoint. Equation \eqref{eq:03-forest:2-2} forces the identical nesting and disjointness relations on any pairwise incident family \((X_e)\). Theorem UB.Laminar, with no prescribed point map, gives a single permutation sending every selected \(A_e\) to \(X_e\). This proves (G3) for every partial flag.

Repeated use of (1) gives \(H_J=K_{E\setminus J}\). Different edge sets give different component partitions and different groups, since every block is nonempty and adding an edge allows a transposition between its two sides. Thus (G1) holds. If at least two edge types remain, the groups obtained by deleting each remaining edge collectively contain \(K_\varnothing\) and each remaining one-edge group, proving (G2). The coset criterion gives the geometry assertions. Its residue group and its subgroups factor over the forest components; all elements from different components are incident, which gives the stated join and the kernel convention. ∎

For unit blocks \(B_v=\{v\}\), write \(\tau_e\) for the endpoint transposition of \(e\). The theorem says \begin{equation}
K_F=\langle\tau_e:e\in F\rangle
 =\prod_C S_C,\qquad K_\varnothing=1.
\label{eq:03-forest:2-3}
\end{equation} This is the unweighted model used for {[}BC98{]}. The weighted formulation also supplies the parity construction after an edge contraction and the disjoint-subset construction after a star specialization.

\section{Overgroups of a two-block stabilizer}
\textbf{Lemma UB.Overgroups (two-block overgroups).} Let \(\Omega=A\sqcup C\) be a finite set, with \(a=|A|\), \(c=|C|\) positive, and put \(H=S_A\times S_C\). If \(a\ne c\), or \(a=c=1\), then \(H\) is maximal in \(S_\Omega\). If \(a=c\geq2\), its only proper strict overgroup is the unordered two-block stabilizer \[
W=(S_A\times S_C)\rtimes C_2,
\] which is maximal and satisfies \([W:H]=2\). In all cases the subgroup interval \([H,S_\Omega]\) is a chain. A positive weighted tree has at most one edge whose two sides have equal total weight.

\emph{Proof.} The case \(a=c=1\) is immediate. The group \(H\) contains the transpositions within each of its two blocks. If an overgroup element does not preserve their unordered partition, its conjugates of these transpositions include one crossing the original blocks. Together with the original within-block transpositions the support graph is connected, so these transpositions generate \(S_\Omega\). This remains valid when one block is a singleton: a permutation moving that singleton produces a cross-block conjugate from the other block. Thus a proper overgroup must preserve the unordered partition. That stabilizer equals \(H\) when the sizes differ and has index two over it when they agree. Its only further overgroup is \(S_\Omega\), by the same argument.

If an edge bisects the total weight, the side of any other edge away from it is a proper positive-weight subset of one of its halves. It has strictly less than half the total weight and cannot be a bisecting side. ∎

\section{Maximum irredundant length and a normal module}
A generating family is irredundant if deleting any member leaves a
proper subgroup. Write $\mu(G)$ for the largest size of an irredundant
generating family of a finite group $G$, and $\mu'(G)$ for the largest
size of any independent family in it. These are maximum lengths,
not the minimum number of generators. For an abelian normal subgroup
$V\unlhd G$, let $\mu'_G(V)$ denote the largest size of a family
independent under generation by all $G$-conjugates.

\begin{lemma}[UB.NormalQuotient]\label{lem:quotient}
For finite $G$ and $N\unlhd G$,
$\mu(G)\leq\mu(G/N)+\mu'(N)$. If $N$ is abelian, the final term
can be replaced by $\mu'_G(N)$. In particular, if $V\unlhd G$ is
an elementary abelian minimal normal subgroup, then
\begin{equation}\label{eq:normalmodule}
 \mu(G)\leq\mu(G/V)+1.                                      \end{equation}
No splitting of the extension is needed.
\end{lemma}
\begin{proof}
Take an irredundant generating family $S$ of $G$, and an inclusion-minimal
subfamily $T$ whose images generate $G/N$. For $s\in S\setminus T$,
choose a word $w_s$ in $T$ with $sN=w_sN$, and put
$v_s=sw_s^{-1}\in N$. If one $v_s$ lies in the subgroup generated by
the others, substitution expresses $s$ using $S\setminus\{s\}$,
a contradiction. Since $|T|\leq\mu(G/N)$, the first bound follows.

When $N$ is abelian, $G=N\gen{T}$ implies that its conjugation action
on $N$ is already the action of $\gen{T}$. Thus an expression for
$v_s$ in conjugates of the other $v_t$ again substitutes to an
expression for $s$ without $s$. This proves the normal-module bound.
If $V$ is elementary abelian minimal normal, it is an irreducible
$\mathbb F_pG$-module. Any nonzero vector generates all of $V$ under
conjugation, so an independent family for this action has at most one
member. This proves \eqref{eq:normalmodule} for maximum irredundant length.
\end{proof}

This is Cameron--Cara's normal-quotient theorem [CC, Theorem 1.1]. Both the extremal and the high-rank applications use this one complete proof. We retain Whiston's bound that an independent family in $S_m$
has at most $m-1$ members, with equality only when it generates $S_m$.
In particular $\mu'(A_m)\leq m-2$. No classification of the C-groups
of $\mathrm{PSL}_2$ or $\mathrm{PGL}_2$ is used below.

\begin{corollary}\label{cor:affine}
If $G$ has a faithful primitive affine action with elementary abelian
translation subgroup $V$, then \eqref{eq:normalmodule} applies. In particular
$\mu(\mathrm{AGL}(4,2))\leq7$.
\end{corollary}
\begin{proof}
The translation module is irreducible: a proper invariant subspace
would give a nontrivial block system. Thus $V$ is minimal normal.
For $\mathrm{AGL}(4,2)$ the quotient is
$\mathrm{GL}(4,2)\cong A_8$, so Whiston gives $\mu(A_8)\leq6$.
The classical exceptional isomorphism is a retained identification;
the finite certificate also checks an explicit isomorphism for the
catalogue representation used below.
\end{proof}

\section{All transitive imprimitive groups of degree nine}
\begin{theorem}\label{thm:nine}
Every transitive subgroup of $S_3\wr S_3$ in its imprimitive action
on nine letters has $\mu(G)\leq5$. Consequently it cannot be a
maximal parabolic of a rank-seven C-group marking of $S_9$.
\end{theorem}
\begin{proof}
Let $K$ be the kernel of the action on the three blocks, and put
$A=K\cap A_3^3$. Then $A$ is an elementary abelian normal 3-subgroup,
viewed as an invariant subspace of $\mathbb F_3^3$. The quotient
$G/A$ has normal subgroup $W=K/A\leq C_2^3$ and transitive quotient
$B\leq S_3$, so $B=C_3$ or $S_3$.

Choose an element of $G$ inducing a three-cycle on the blocks.
On $A_3^3$ its conjugation action is a cyclic permutation of the
coordinates with coefficients $\varepsilon_i\in\{1,-1\}$.
The vector in any one coordinate is cyclic for this operator, whose
characteristic and minimal polynomials are
$X^3-\delta=(X-\delta)^3$, where $\delta=\varepsilon_1\varepsilon_2
\varepsilon_3\in\mathbb F_3$. Thus the full module is
$\mathbb F_3[X]/(X-\delta)^3$. Its invariant subspaces form a chain,
and each is cyclic. Any family of vectors in an invariant subspace
therefore has one vector whose cyclic span contains all the others.
Independence under the whole group implies independence under this
one element, so $\mu'_G(A)\leq1$.

The corresponding action on $C_2^3$ is ordinary coordinate permutation.
Under the three-cycle it is the direct sum of a trivial line and the
irreducible plane with polynomial $X^2+X+1$. Every invariant subspace
is a sum of a subset of these two nonisomorphic irreducibles. The two
polynomial projection operators show that a vector generates each
component on which it is nonzero. At most two vectors are therefore
needed from any given family to generate its cyclic submodule; a
family independent under the full group has at most two members.
Hence $\mu'_{G/A}(W)\leq2$.

Twice applying Lemma~\ref{lem:quotient} gives
\[
 \mu(G)\leq\mu(G/A)+1\leq\mu(B)+2+1\leq5.
\]
This applies to the maximum irredundant length, not just a selected
small generating set. A parabolic of rank six in a C-group has six
independent distinguished generators, contradicting the bound.
\end{proof}

\section{Independent vectors in a transitive binary module}
The following elementary observation is useful beyond the present
classification. It bounds irredundance of a particular family, rather
than merely providing some short generating family for the module.

\begin{lemma}\label{lem:simplequotients}
Let $M$ be a finite-dimensional module over a finite field, generated
irredundantly as a module by $v_1,\ldots,v_s$. There are simple quotients
$L_1,\ldots,L_s$ of $M$ and a surjection
$M\longrightarrow\bigoplus_{i=1}^s L_i$.
\end{lemma}
\begin{proof}
For each $i$, choose a maximal proper submodule $M_i$ containing all
$v_j$ with $j\ne i$. It does not contain $v_i$, since those vectors
generate $M$. Put $L_i=M/M_i$. Under the product quotient map, the
submodule generated by $v_i$ has zero coordinates except the $i$th,
where its image is the nonzero submodule $L_i$. Thus the image contains
every direct summand and is the whole direct sum.
\end{proof}

\begin{lemma}\label{lem:binary}
Let $P\leq S_m$ be transitive, where $m\geq4$, acting by coordinate
permutation on $V=\mathbb F_2^m$. For every invariant subspace
$W\leq V$, every $P$-independent family in $W$ has size at most $m-2$.
For $P=S_m$ the stronger bound two holds.
\end{lemma}
\begin{proof}
Let $M$ be the submodule generated by an independent family of size $s$.
If $M<V$ and $s\geq m-1$, Lemma~\ref{lem:simplequotients} forces
$\dim M=s=m-1$ and every simple quotient there has dimension one.
Over $\mathbb F_2$ a one-dimensional representation is trivial, so the
surjection is an isomorphism and $M$ is fixed pointwise by $P$.
But the fixed space of a transitive permutation group on $V$ is the
one-dimensional constant subspace, a contradiction.

If $M=V$, at most one of those simple quotients has dimension one:
otherwise the product surjection would give two linearly independent
maps from $V$ to the trivial module, whereas a $P$-invariant linear
functional has equal coefficients in all $m$ positions. Hence
$m\geq2s-1$, and $s\leq\lfloor(m+1)/2\rfloor\leq m-2$.

For the final assertion, if every vector in the family is constant,
one suffices. Otherwise choose a nonconstant vector $v$. Subtracting
its image under a transposition that exchanges a zero and a one gives
$e_i+e_j$. Conjugation then generates the entire even-weight subspace.
This subspace has codimension one, so at most one further member of the
family is needed. Independence therefore permits at most two.
\end{proof}

\begin{theorem}\label{thm:twoblocks}
Every transitive subgroup of $S_2\wr S_m$ in degree $2m$, $m\geq5$,
has maximum irredundant generating length at most $2m-4$.
\end{theorem}
\begin{proof}
Let $W\leq\mathbb F_2^m$ be the block kernel and $P\leq S_m$ the
transitive block image. The conjugation action on $W$ is coordinate
permutation. If $\mu(P)\leq m-2$, Lemmas~\ref{lem:quotient}
and~\ref{lem:binary} give $\mu(G)\leq2m-4$. Otherwise Whiston's
equality clause forces $P=S_m$. The sharper module bound gives
$\mu(G)\leq(m-1)+2=m+1\leq2m-4$.
\end{proof}

\begin{lemma}\label{lem:smallbase}
For $m\geq4$, every transitive subgroup of $S_3\wr S_m$ has
$\mu(G)\leq3m-4$. For $m\geq3$, every transitive subgroup of
$S_4\wr S_m$ has $\mu(G)\leq4m-4$.
\end{lemma}
\begin{proof}
For blocks of size three, write $K$ for the block kernel,
$A=K\cap A_3^m$, $W=K/A\leq\mathbb F_2^m$, and $P=G/K$.
If $P\ne S_m$, then $\mu(P)\leq m-2$, the binary-module cost
is at most $m-2$, and the ternary-module cost is at most its
ambient dimension $m$. This gives $3m-4$.

If $P=S_m$, the binary cost is at most two. The ternary cost is at
most $m-1$. Indeed, an independent family of $m$ vectors would span
all of $\mathbb F_3^m$, and Lemma~\ref{lem:simplequotients} would
identify this monomial module with a sum of one-dimensional modules.
Its action would then have abelian image. But its permutation image
on the coordinate axes is $S_m$, which is nonabelian. Thus the total
cost is at most $(m-1)+2+(m-1)=2m\leq3m-4$.
For three blocks of size three use Theorem~\ref{thm:nine} instead.

Now let the blocks have size four. The block-kernel projections are
conjugate. If they are proper subgroups of $S_4$, every independent
family in each projection has at most two members, by Whiston's
equality clause. Successive projections of the kernel and
Lemma~\ref{lem:quotient} bound its independent families by $2m$.
Consequently $\mu(G)\leq(m-1)+2m\leq4m-4$.

Otherwise every projection is $S_4$. Put $A=K\cap V_4^m$, where
$V_4$ is the normal Klein four subgroup of $S_4$. Then $G/A$ has a
faithful transitive imprimitive action with $m$ blocks of size three,
so the preceding bound gives $\mu(G/A)\leq3m-4$.
We claim that the normal-module cost of $A\leq\mathbb F_2^{2m}$
is at most $m$. The group $L=K/A$ is a subdirect subgroup of $S_3^m$.
Its normal subgroup $Q=L\cap C_3^m$ projects onto $C_3$ in every
coordinate: the commutator subgroup of $L$ already has this property.
Thus $Q$ fixes no nonzero vector in the full $2m$-dimensional module,
since the natural $C_3$ action on each two-dimensional coordinate
has no fixed vector. Averaging over the odd-order group $Q$ shows
that taking $Q$-fixed vectors is exact, so every subquotient of this
module also has zero $Q$-fixed space. In particular every simple
$G$-module quotient in Lemma~\ref{lem:simplequotients} has dimension
at least two. Hence a normal-independent family in $A$ has size
at most $\dim A/2\leq m$. The normal-module inequality now gives
$\mu(G)\leq(3m-4)+m=4m-4$.
\end{proof}

\section{A bounded primitive certificate}
For an integer $a>0$, let $\Omega(a)$ be its number of prime factors
with multiplicity. Every strictly ascending subgroup chain in a finite
group of order $a$ has at most $\Omega(a)$ strict inclusions, because
the product of its successive integral indices divides $a$. Thus the
C-rank is at most $\Omega(|G|)$.

The required program \texttt{primitive\_bound.g} enumerates the complete
GAP primitive-group catalogue in degrees 9--25, omitting $A_n,S_n$,
and reports exactly twelve groups for which this elementary bound
exceeds $n-4$. All remaining catalogue entries are thereby discharged.
The twelve exceptions are:
\begin{equation}\label{eq:primitiveexceptions}
\begin{array}{c|l}
9&3^2\!:\!QD_{16},\quad3^2\!:\!SL_2(3),\quad AGL_2(3),\quad
 PSL_2(8),\quad PSL_2(8)\!:\!3;\\
10&A_6\!:\!2,\quad S_6,\quad A_6\!.\!2,\quad A_6\!:\!2^2;\\
11&M_{11};\\12&M_{12};\\16&AGL_4(2).
\end{array}                                                   \end{equation}
The program verifies the precise subgroup and quotient identifications
used next. Its retained inputs are the completeness of this catalogue
and GAP's exact finite-group algorithms; it performs no search for
C-group markings and imports none of the source's reported C-ranks.

For the first two degree-nine groups the irreducible translation module
costs one and the quotient has at most four prime factors, giving five.
For $AGL_2(3)$ the quotient $GL_2(3)$ has Frattini subgroup of order two
and Frattini quotient $S_4$. Passing to a Frattini quotient preserves
irredundant generating families: a subset generates the whole group
exactly when its image generates the quotient. Thus $\mu(GL_2(3))
=\mu(S_4)\leq3$, and the affine group has $\mu\leq4$.
Every maximal subgroup of $PSL_2(8)$ has order 56, 18 or 14; these have
at most four prime factors, so a C-group marking has rank at most five.
The extension by 3 is not generated by its involutions: the program
checks that they generate only the index-three subgroup.

In degree ten the three groups other than $S_6$ have normal socle
$A_6$ and quotient of order two or four. Lemma~\ref{lem:quotient}
and $\mu'(A_6)\leq4$ give at most six generators; $S_6$ has at most
five. Corollary~\ref{cor:affine} treats degree sixteen.

For $M_{11}$ the maximal-subgroup orders are
$720,660,144,120,48$. The subgroup generated by all involutions in
each has at most six prime factors; in the order-720 case its order
is 360. Any C-group parabolic lying in such a maximal subgroup lies
in its involution-generated subgroup. Hence $M_{11}$ has C-rank at
most seven. For $M_{12}$ the maximal-subgroup orders are
\[
7920,7920,1440,1440,660,432,432,240,192,192,72.
\]
The first two are $M_{11}$. The order-1440 groups have normal $A_6$
with quotient order four, so their irredundant lengths are at most six.
Every remaining maximal subgroup has at most seven prime factors.
Deleting one distinguished generator and placing the resulting proper
parabolic in a maximal subgroup gives C-rank at most eight for $M_{12}$.
All these finite assertions are checked by the complete program.

For arbitrary degree $n\geq26$, retain Mar\'oti's Theorem 1.1 in its
original form, including the lower containment $A_m^\ell\leq G$ in
the product-action alternative. If $\ell\geq2$, the same abstract
group acts faithfully on $m\ell$ points, so Whiston bounds its
independent families by $m\ell-1\leq m^\ell-4\leq n-4$; the
inequality holds for $m\geq2$, $\ell\geq2$, $m^\ell\geq9$.
If $\ell=1$, the natural action case is excluded by $G\ne A_n,S_n$;
otherwise the $k$-subset action has $k\geq2$ and $m\leq n/2$,
so the same conclusion follows from $m-1\leq n-4$.
The Mathieu alternative has degree at most 24.

In the remaining alternative, Mar\'oti bounds the group order by
$M(n)=n\prod_{i=0}^{\lfloor\log_2n\rfloor-1}(n-2^i)$.
For $26\leq n\leq63$, exact integer arithmetic in the certificate
verifies $M(n)<2^{n-3}$. For $n\geq64$, put
$k=\lfloor\log_2n\rfloor\geq6$. Then
\[
 \log_2M(n)<(k+1)\log_2n<(k+1)^2\leq2^k-3\leq n-3.
\]
The middle inequality follows at $k=6$ and is preserved when $k$
increases by one. Since C-rank is an integer and is less than $n-3$,
it is at most $n-4$. We have proved the full primitive-parabolic
bound for all degrees, conditional only on the explicitly retained
Whiston and Mar\'oti theorems and the bounded exact catalogue inputs.

\section{A uniform simple-kernel bound}
\begin{lemma}\label{lem:scott}
Let $S$ be a finite nonabelian simple group. Every subdirect subgroup
of $S^m$ is a product of diagonal copies of $S$ on a partition of the
$m$ coordinates, with automorphisms allowed in the diagonal maps.
A strictly increasing chain of subdirect subgroups has at most $m-1$
strict steps.
\end{lemma}
\begin{proof}
Proceed by induction on $m$. Project onto the first $m-1$ coordinates.
The kernel in the last factor is normal in $S$, so is either $S$ or
trivial. In the first case the subgroup is the product of its projection
with $S$. In the second it is the graph of an onto homomorphism from
that projection to $S$. By induction the projection is a direct product
of simple diagonal factors. Their images are commuting normal subgroups
of $S$, so exactly one has nontrivial image, and its map is an
isomorphism. The last coordinate therefore joins that diagonal block.
This proves the description. In an inclusion of such subgroups the
larger group can only refine the partition. If the number of blocks
does not increase, their orders are equal and the groups coincide.
There are between one and $m$ blocks, proving the chain bound.
\end{proof}
This is the usual subdirect-product lemma, often called Scott's lemma;
its complete elementary proof is included here, rather than retained
as an unpriced extra premise.

\begin{lemma}\label{lem:simplekernel}
Suppose $N\unlhd G$ is a subdirect subgroup of $S^m$ for a finite
nonabelian simple group $S$, and conjugation by $G$ permutes the
coordinates transitively, using automorphisms on individual factors.
Then
\begin{equation}\label{eq:simplekernel}
 \mu(G)\leq\mu(G/N)+\mu'(S)+m-1.                    \end{equation}
\end{lemma}
\begin{proof}
Normalize an irredundant generating family as in
Lemma~\ref{lem:quotient}, obtaining $T\cup E$, with $T$ minimally
surjecting onto $G/N$ and $E\subset N$. The normalized extra elements stay independent relative to $\gen{T}$,
by the substitution argument, and no member of $T$ can be dropped since
its quotient image is indispensable. Thus this replacement preserves
irredundance. Write $T_0=\gen{T}$ and $N_0=T_0\cap N$.
The subgroup generated by all $T_0$-conjugates of $E$ is normalized by
$T_0$ and contains $E$, hence is normal in $G$. It follows that
$N$ is generated by $N_0$ and these conjugates.

Choose a minimal subset $B\subseteq E$ such that the projections of
$N_0$ and of all $T_0$-conjugates of $B$ generate $S$ in the first
coordinate. Then $|B|\leq\mu'(S)$. To see this, for each $b\in B$
choose an element in the projected subgroup generated by its conjugates
which does not belong to the subgroup generated by $N_0$ and the
corresponding subgroups for the other members of $B$. These selected
elements form an independent family in $S$.

Put $J=\gen{T,B}$. The group $J\cap N$ projects onto the first copy
of $S$. Conjugation by $N\leq S^m$ fixes each coordinate, so the transitive
coordinate action of $G$ is already that of $T_0$. Therefore
it projects onto every copy. Add the elements of $E\setminus B$
one by one. Each addition strictly enlarges the generated group, by
irredundance. Since the image modulo $N$ stays fixed, it strictly
enlarges its intersection with $N$. These intersections form a chain
of subdirect subgroups, with at most $m-1$ strict steps by
Lemma~\ref{lem:scott}. Thus
$|T|+|E|\leq\mu(G/N)+\mu'(S)+m-1$, as required.
\end{proof}

\section{Primitive block stabilizers}
\begin{lemma}\label{lem:primitiveindependent}
Let $L$ be primitive of degree $k\geq5$, with $L\ne A_k,S_k$.
Then $\mu'(L)\leq k-3$, except for the degree-six copy of $S_5$,
which has $\mu'=4$.
\end{lemma}
\begin{proof}
For $k\geq26$, the Mar\'oti and Whiston argument above bounds all
independent families, not only C-group markings, by $k-4$.
For $5\leq k\leq25$, the complete primitive catalogue and the
prime-factor bound leave exactly the following eleven entries:
\begin{equation}\label{eq:independentexceptions}
\begin{array}{c|l}
5& C_5:C_4\\
6& A_5,\ S_5\\
7& PSL_3(2)\\
8& AGL_3(2),\ PGL_2(7)\\
9& AGL_2(3),\ PSL_2(8):3\\
10& A_6:2^2\\
12& M_{12}\\
16& AGL_4(2).
\end{array}                                                    \end{equation}
The same required program checks this complete list and the exact
maximal-subgroup, normal-subgroup and quotient facts below.

The group $C_5:C_4$ has irredundant length at most two by its minimal
normal $C_5$ and cyclic prime-power quotient. Every proper subgroup
is cyclic or dihedral of order ten, hence also has bound two.
Whiston treats $A_5$ and $S_5$ in their natural degrees five.
Every maximal subgroup of $PSL_3(2)$ is $S_4$ or has at most two prime
factors, so a generating family has length at most four and a proper
subgroup has independent length at most three. The index-two extension
$PGL_2(7)$ consequently has independent length at most five.
Here and below we use the immediate subgroup version of
Lemma~\ref{lem:quotient}: $\mu'(G)\leq\mu'(G/N)+\mu'(N)$.

For $AGL_3(2)$ consider any subgroup $H$, its intersection $U$ with
the translation module $V$, and its linear image $P\leq GL_3(2)$.
If $\dim U=0$, the bound is four. If $\dim U=1$ or two, $P$ lies
in a line or plane stabilizer, each isomorphic to $S_4$ (its normal $2^2$ and quotient $GL_2(2)$
act as the full affine group on four points), so the bound
is $3+\dim U\leq5$. If $U=V$ and $P$ is irreducible, the normal
module costs one and the quotient costs at most four. If $P$ is
nontrivial and reducible, it lies in an $S_4$ stabilizer, and the
module costs at most two: three independent vectors would, by
Lemma~\ref{lem:simplequotients}, identify $V$ with a sum of three
one-dimensional, hence trivial, modules, contradicting the faithful
nontrivial action. Finally $P=1$ costs only three. Thus every $H$
has irredundant length at most five.

Every proper subgroup of $GL_2(3)$ lies in a maximal subgroup whose
order has at most four prime factors; the whole group has irredundant
length three by its Frattini quotient $S_4$. Thus $\mu'(GL_2(3))\leq4$,
and $\mu'(AGL_2(3))\leq2+4=6$. The earlier maximal-subgroup bound
for $PSL_2(8)$ gives independent length at most five, so its index-three
extension has bound six. The degree-ten group has normal $A_6$ and
quotient of order four, giving at most $4+2=6$. Every maximal subgroup
of $M_{12}$ has at most eight prime factors, so its independent length
is at most nine. Finally $AGL_4(2)$ has a four-dimensional normal
translation module and quotient $A_8$, giving $\mu'\leq4+6=10$.
All satisfy the asserted bounds.
\end{proof}

\section{The full imprimitive bound}
\begin{theorem}\label{thm:imprimitive}
Every transitive imprimitive permutation group of degree $n\geq9$
has maximum irredundant generating length at most $n-4$.
\end{theorem}
\begin{proof}
Choose a minimal nontrivial block, of size $k$, and let $m\geq2$
be the number of its translates, so $n=km$. The induced group $L$
of the setwise block stabilizer is primitive: a proper block for $L$
would propagate to a smaller block system for $G$. After choosing
identifications of the blocks by elements of $G$, the group embeds
in $L\wr S_m$. Its block image $P$ is transitive, and the projection
$K_1$ of the block kernel $K$ is normal in $L$; all other projections
are conjugate to it.

Blocks of size two, three or four are covered by
Theorems~\ref{thm:nine}, \ref{thm:twoblocks} and
Lemma~\ref{lem:smallbase}. Assume $k\geq5$.
If $m\geq3$ and $L<S_k$, then every subgroup of $L$ has independent
length at most $k-2$, by Whiston's equality clause. Successive
coordinate projections give $\mu'(K)\leq m(k-2)$, and therefore
\[
 \mu(G)\leq(m-1)+m(k-2)=n-m-1\leq n-4.
\]
If $m=2$ and $L\ne A_k,S_k$ is not the degree-six copy of $S_5$,
Lemma~\ref{lem:primitiveindependent} similarly gives
$\mu(G)\leq1+2(k-3)=n-5$.

It remains to treat $L=A_k,S_k$, and the exceptional degree-six
copy of $S_5$ with $m=2$. First take $L=A_k$ or $S_k$.
By simplicity of $A_k$ and its trivial centralizer in $S_k$, the normal subgroup $K_1$ is either trivial,
$A_k$, or $S_k$. In the trivial case $K=1$ and
$\mu(G)\leq m-1\leq n-4$. Otherwise
$N=K\cap A_k^m$ is subdirect in $A_k^m$: the commutator subgroup
of $K$ already projects onto $A_k$ in every coordinate. The quotient
$G/N$ embeds in $C_2\wr S_m$. Its maximum irredundant length is at
most $(m-1)+m=2m-1$ by the normal-module bound, so
Lemma~\ref{lem:simplekernel} and $\mu'(A_k)\leq k-2$ give
\[
 \mu(G)\leq (2m-1)+(k-2)+(m-1)=k+3m-4.
\]
For $k\geq5,m\geq3$ this is at most $km-4$, since
$(k-3)(m-1)\geq3$.
When $m=2$, the quotient lies in $C_2\wr S_2\cong D_8$;
every subgroup of $D_8$ has irredundant length at most two.
The sharper bound is then $\mu(G)\leq2+(k-2)+1=k+1\leq2k-4$.

Finally let $k=6,m=2,L\cong S_5$. Use the characteristic subgroup
$A_5\lhd L$ and the same argument in the abstract factors.
The block kernel is trivial or projects onto $A_5$ or $S_5$.
In the latter two cases Lemma~\ref{lem:simplekernel}, now with
$S=A_5$, gives $\mu(G)\leq2+3+1=6\leq12-4$.
This completes all cases and every parameter value.
\end{proof}

\begin{corollary}\label{cor:intransitive}
Let $(S_n;\rho_0,\ldots,\rho_{r-1})$ be a C-group with
$n\geq9$ and $r\in\{n-2,n-1\}$. Every maximal standard
parabolic is intransitive on the $n$ letters.
\end{corollary}
\begin{proof}
Each such parabolic has an independent distinguished family of
$r-1\geq n-3$ elements. The primitive bound and
Theorem~\ref{thm:imprimitive} exclude every transitive proper
subgroup except $A_n$. Suppose one parabolic is $A_n$ and let
$a$ be the omitted involution. For each other distinguished
involution $b$, the intersection property says
$A_n\cap\gen{a,b}=\gen{b}$. Since $aba\in A_n$, this forces
$aba=b$. Thus $a$ centralizes all of $A_n$ and
$S_n=A_n\times\gen{a}$, impossible. The centralizer of $A_n$
in $S_n$ is trivial for $n\geq4$, as is seen by commuting with
all three-cycles. Hence no maximal parabolic is transitive.
\end{proof}


\section{The retained extremal inputs}
The classification proved in this section is Cameron--Cara, Theorem 2.1, with its strong intersection and relative-subgroup consequences. Its proof is now common to two applications. The ordinary and parity-modified families remain their source results; the present organization makes no new historical classification claim.
\textbf{W1 (Whiston bound and extremal generation).} An independent family in \(S_m\) has at most \(m-1\) members, and an independent family with \(m-1\) members generates \(S_m\). We use precisely this implication, not a converse asserting that every independent generating family has maximum size.

\textbf{W2 (Whiston omission consequence).} If \(n\geq7\) and \((s_i)\) is an independent generating family of \(S_n\) of size \(n-1\), each omission subgroup \(\langle s_j:j\ne i\rangle\) is intransitive, is transitive imprimitive with blocks of size two, or is \(A_n\). This is the separate consequence of Whiston's proof stated on page 4 of the 2002 manuscript. Its dependence on the classification of finite simple groups is retained. W2 is not deduced here from W1.
Cayley's labelled-tree count and the innerness of every automorphism of \(S_n\) for \(n\geq7\) are retained classical inputs. The CFSG-dependent omission consequence W2 is separate from W1.

\section{Tree markings and the parity construction}
\textbf{Theorem UB.Sign (sign-kernel extension).} Let the blocks be singletons and distinguish an edge \(d\), with \(s=\tau_d\). For \(e\ne d\), choose \(\epsilon_e\in\{1,-1\}\) and put \begin{equation}
r_d=s,\qquad r_e=(s\tau_e)^{\epsilon_e}.
\label{eq:maximal-symmetric-markings:2-4}
\end{equation} For \(U\subseteq E\), set \(F(U)=\{d\}\cup(U\setminus\{d\})\). Then \begin{equation}
\langle r_e:e\in U\rangle=
\begin{cases}
K_{F(U)},&d\in U,\\
K_{F(U)}\cap A_\Omega,&d\notin U.
\end{cases}
\label{eq:maximal-symmetric-markings:2-5}
\end{equation} Both the ordinary marking \((\tau_e)\) and \eqref{eq:maximal-symmetric-markings:2-4} generate \(S_\Omega\), are independent, and satisfy the full subset intersection property with \(B=1\).

\emph{Proof.} Inverting an individual generator leaves its cyclic subgroup unchanged. If \(s\) is included, \(s\) and \(s\tau_e\) recover \(\tau_e\), proving the first case. The sign kernel of \(K_F\), for \(d\in F\), is generated by the \(s\tau_e\), \(e\in F\setminus\{d\}\): every even word in edge transpositions is a product of pairs, and \[
\tau_e\tau_f=(s\tau_e)^{-1}(s\tau_f),\qquad s\tau_d=1.
\] This proves the second case, including the empty set, since \(K_{\{d\}}\cap A_\Omega=1\). Intersect \eqref{eq:maximal-symmetric-markings:2-5} using Theorem UB.F and the sign condition. Since \(F(U)\cap F(V)=F(U\cap V)\), the result is exactly the subgroup for \(U\cap V\). The same conclusion for ordinary markings follows from the unit-weight specialization of Theorem UB.F. Omitting an ordinary edge or a twisted label \(e\ne d\) leaves its cut invariant; the omitted generator crosses that cut. Omitting \(r_d\) leaves only even permutations, whereas \(r_d\) is odd. These observations prove independence and the generation assertions. ∎
\section{The extremal converse}
\textbf{Corollary UB.Blocks (the two-point-block obstruction).}
If $m\geq4$ and $H\leq C_2\wr S_m$ is transitive on $2m$ points,
then $H$ has no independent generating family of size $2m-2$.
\begin{proof}
Let $N$ be the block kernel and $P$ the transitive block image.
If $P<S_m$, Whiston's equality clause gives $\mu(P)\leq m-2$,
and Lemmas~\ref{lem:quotient} and~\ref{lem:binary} give
$\mu(H)\leq2m-4$. If $P=S_m$, their stronger module bound gives
$\mu(H)\leq m+1$. Both quantities are strictly smaller than $2m-2$
for $m\geq4$, including the endpoint $m=4$.
\end{proof}

\textbf{Theorem UB.Extremal (extremal recognition).} Assume \(n\geq7\). Every independent family of \(n-1\) elements of \(S_n\) is an ordinary unit-weight tree marking from Theorem UB.F or a sign-kernel marking \eqref{eq:maximal-symmetric-markings:2-4}. If the family is independent relative to \(B\leq S_n\), then \(B=1\).

\emph{Proof.} By W1 the family generates \(S_n\). For its omission groups \(G_i\), W2 and Corollary UB.Blocks exclude the two-point-block case: in that case \(n=2m\geq8\), and the remaining \(n-2\) generators are an independent generating family for \(G_i\). At most one \(G_i\) is \(A_n\), because its omitted generator must be the unique odd member. Label a possible exception by \(1\).

For each \(i>1\), choose an edge \(e_i=\{x,s_i(x)\}\) crossing two \(G_i\)-orbits. Such an edge exists, since otherwise every generator would preserve those orbits. Choose \(e_1\) from any moved point of \(s_1\). If all the omission groups are intransitive, choose \(e_1\) crossing its own omission orbits as well. For \(j\ne i\), the edge \(e_j\) lies inside a \(G_i\)-orbit because \(s_j\in G_i\). In particular the chosen edges are distinct: a repeated edge with a label \(i>1\) would both cross and stay inside the \(G_i\)-orbits. Thus a cycle cannot contain \(e_i\) for \(i>1\): the remaining edges would give a path within a \(G_i\)-orbit between its endpoints. A cycle cannot consist of \(e_1\) alone. The \(n-1\) edges therefore form a tree on the \(n\) points.

Deleting \(e_i\), for \(i>1\), leaves two connected sets each contained in a \(G_i\)-orbit; intransitivity makes them exactly its two orbits. Hence \(s_j\) preserves this cut whenever \(j\ne i\). Intersect these cut stabilizers using Theorem UB.F. It follows that \(s_1\in K_{\{e_1\}}\), so \(s_1=\tau_{e_1}\), and that \(s_i\in K_{\{e_1,e_i\}}\). If every omission group is intransitive, include the first cut too to obtain \(s_i=\tau_{e_i}\) for every \(i\). Otherwise \(G_1=A_n\), and all other generators are nonidentity even elements of the two-edge group. For adjacent edges that group is \(S_3\), whose nonidentity even elements are the two inverse three-cycles; for disjoint edges it is \(C_2\times C_2\), with a unique nonidentity even element. These are precisely \((s_1\tau_{e_i})^{\pm1}\), giving \eqref{eq:maximal-symmetric-markings:2-4}.

Now suppose the original family is independent relative to \(B\). It is independent in the ordinary sense, so the preceding classification applies. By Theorems UB.F and UB.Sign its omission groups are the tree-cut stabilizers, except that in the sign-kernel case the distinguished omission group is \(A_n\). At each maximal omission group independence forces \(B\leq G_i\); otherwise \(\langle B,G_i\rangle=S_n\). A pendant edge different from the distinguished edge exists and yields a point stabilizer, so \(B\) fixes a point. At a nonmaximal cut, Lemma UB.Overgroups forces \(B\) into the unordered two-block stabilizer. Its elements outside the ordered stabilizer interchange the halves and fix no point. Thus \(B\leq G_i\) here also. Strong independence from Theorems UB.F and UB.Sign gives \(B\leq\bigcap_iG_i=1\). ∎


\section{A finite recognition interface}
For a partition $\pi$ of a finite set $Q$, write
$\Y(\pi)=\prod_{B\in\pi}\Sym(B)$, as a subgroup of $\Sym(Q)$.
Let $(h_i)_{i\in I}$ be a finite family in $\Sym(Q)$ and let
$H_S=\langle h_i:i\in S\rangle$ for $S\subseteq I$.
Fix $d\in Q$. Suppose each $S$ is assigned a partition $\pi_S$ of $Q$.
If $x\notin Q$, let $\pi_S^x$ be the partition of $Q\cup\{x\}$
obtained by adjoining $x$ to the block containing $d$.

We impose the following finite conditions, for every $S,R\subseteq I$:
\begin{align}
H_S\cap H_R&=H_{S\cap R}, &
H_S\cap\Y(\pi_R)&=H_{S\cap R},\label{eq:finite1}\\
\pi_S\wedge\pi_R&=\pi_{S\cap R}, &
\langle H_S,(d\ x)\rangle&=\Y(\pi_S^x).
\label{eq:finite2}
\end{align}
Here the meet is common refinement. In particular, taking $S=R$ in
\eqref{eq:finite1} gives $H_S\leq\Y(\pi_S)$. These hypotheses concern
subgroups on $|Q|$ or $|Q|+1$ letters; they do not assume the desired
intersection property on any larger degree.

Let $T$ be a finite tree meeting $Q$ in exactly $\{d\}$, with at least
one edge, and suppose $d$ is a leaf of $T$. Put
$\Omega=Q\cup V(T)$. Alongside the $h_i$, take the transposition
$t_e$ exchanging the endpoints of each edge $e$ of $T$. All these
permutations fix the unspecified letters.

\begin{theorem}\label{thm:attachment}
Under \eqref{eq:finite1}--\eqref{eq:finite2}, the family
$(h_i)_{i\in I},(t_e)_{e\in E(T)}$ has the full subset intersection
property. If the $h_i$ are distinct nonidentity involutions, then so is
this enlarged family and it is a C-group marking of its generated group.
If moreover $\pi_I=\{Q\}$, that group is $\Sym(\Omega)$.
\end{theorem}

\begin{proof}
For $F\subseteq E(T)$, let $D_F$ be the vertex set of the component
containing $d$ in $(V(T),F)$; the other component sets will be denoted
$C$. The tree transpositions on a connected component generate its full
symmetric group: along the unique paths, conjugating endpoint
transpositions produces every transposition from one fixed vertex.
Thus
\begin{equation}\label{eq:all-subsets}
\langle h_i\ (i\in S),t_e\ (e\in F)\rangle
= H_S[D_F]\times\prod_{C\ne D_F}\Sym(C),
\end{equation}
where $H_S[\{d\}]=H_S$. If $D_F\ne\{d\}$, then $H_S[D_F]$
is the Young subgroup obtained from $\pi_S$ by adjoining
$D_F\setminus\{d\}$ to its $d$-block. Indeed, choose any
$x\in D_F\setminus\{d\}$. The subgroup $\Sym(D_F)$ contains
$(d\ x)$, so \eqref{eq:finite2} supplies every transposition within
each unchanged core block and within the enlarged $d$-block on
$Q\cup\{x\}$. Conjugation by $\Sym(D_F)$ supplies the same
transpositions for every point of $D_F$. Conversely all the original
generators preserve the claimed blocks. This proves \eqref{eq:all-subsets}.

We now intersect the groups for $(S,F)$ and $(R,E)$. The partition calculation in UB.F applies to their forest components:
common refinement is the component partition of $F\cap E$.
The remaining issue is the core factor, which need not itself be a Young subgroup.

If $D_F=D_E=\{d\}$, the core intersection is $H_{S\cap R}$
by \eqref{eq:finite1}; all other factors follow from the forest
argument. Suppose exactly one of $D_F,D_E$ is a singleton, say
$D_F=\{d\}$. Every permutation in the first group preserves $Q$
setwise and acts there in $H_S$. Its additional membership in the
second group restricts this action exactly to $\Y(\pi_R)$.
Equation \eqref{eq:finite1} gives $H_{S\cap R}$ on $Q$.
Off $Q$, the two forest component partitions have the refinement
just described; their intersections with the enlarged root component
are connected subtrees too. Consequently the remaining factors are
exactly the symmetric groups on the nonroot components of $F\cap E$.

Finally suppose both root components are nonsingletons. The unique
edge incident with $d$ lies in both $F$ and $E$. Therefore
$D_{F\cap E}=D_F\cap D_E$ is also nonsingleton. Both groups in
\eqref{eq:all-subsets} are now Young subgroups of $\Sym(\Omega)$.
Their intersection is the Young subgroup of the common refinement.
On the core this is $\pi_S\wedge\pi_R=\pi_{S\cap R}$;
the root block acquires exactly $D_F\cap D_E\setminus\{d\}$,
and all other blocks are the nonroot components of $F\cap E$.
This is precisely \eqref{eq:all-subsets} for $(S\cap R,F\cap E)$.
All subset intersections have been proved.

No tree transposition is supported on $Q$, and different edges give
different transpositions. Thus the distinctness assertion follows.
For the whole edge set, $D_F=V(T)$ is nonsingleton; if $\pi_I$
has one block, \eqref{eq:all-subsets} is $\Sym(\Omega)$.
\end{proof}


\section{A uniform presentation for a tree}
Let $T$ be a finite tree on $n$ vertices. Write $s_e$ for its edge
transpositions. The following is the presentation in Fernandes--Leemans,
Proposition 3.1, without its unnecessary restriction $n\geq7$.
\begin{theorem}\label{thm:tree}
The symmetric group on the vertices of $T$ has generators $s_e$ and
relations
\[
 s_e^2=1,\qquad (s_es_f)^2=1\quad(e,f\text{ disjoint}),\qquad
 (s_es_f)^3=1\quad(e,f\text{ meet}),
\]
together with $(s_es_fs_es_g)^2=1$ for every three distinct edges
$e,f,g$ incident with one vertex. For each unordered triple it is enough to impose the relation in one
ordering; all the other orderings then follow.
\end{theorem}
\begin{proof}
First observe that the choice of ordering is immaterial. For three
incident involutions $a,f,g$, the pairwise braid relations give
\[
 [afa,g]=1\ \Longleftrightarrow\ [a,fgf]=1
 \ \Longleftrightarrow\ [a,gfg]=1
 \ \Longleftrightarrow\ [aga,f]=1.
\]
The first and last equivalences follow by conjugation and a braid
identity. Also $afa=faf$ interchanges the first two letters.
These interchanges generate all orderings of the triple.

Let $P(T)$ be the presented group. The indicated transpositions satisfy
the relations and generate $\Sym(V(T))$, so there is a surjection
$P(T)\to\Sym(V(T))$. We prove $|P(T)|\leq n!$ by induction on $n$.
The cases $n\leq2$ are immediate. Remove a leaf $p$, with neighbor $q$,
and write $a=s_{pq}$. By induction the subgroup $H$ generated by the
remaining edge generators is a quotient of $\Sym(V(T)\setminus\{p\})$.
Its displayed permutation image is that entire symmetric group, so the
quotient is faithful and we may identify $H$ with $S_{n-1}$.

We first prove that $a$ centralizes $H_q$, the point stabilizer of $q$
in $H$. It commutes with every edge not incident with $q$. If
$f=s_{qu}$ and $g=s_{qv}$ are two distinct other edges at $q$, the
triangle relation says that $afa$ commutes with $g$. The braid
identity $afa=faf$ turns this into the assertion that $a$ commutes
with $fgf=(uv)$. The edge transpositions within the components of
$T-\{p,q\}$, together with these transpositions between their neighbors
at $q$, form a connected transposition graph on all vertices other than
$p,q$. They therefore generate $H_q$. This proves the assertion,
including the case where there is at most one such component.

Choose a neighbor $u$ of $q$ in $T-p$. The relation for $a$ and
$s_{qu}$ is the braid relation. Conjugating it by $H_q$, which acts
transitively on the remaining $n-2$ vertices and centralizes $a$, gives
\[
 af_va=f_vaf_v\qquad\text{for }f_v=(qv),\quad v\ne p,q.
\]
Set $f_q=1$. Every element of $H$ has a unique left coset of $H_q$,
represented by some $f_v$. Consequently the union of the following
at most $n$ left cosets of $H$ is stable under right multiplication by
$H$:
\[
 H,\qquad Ha f_v\quad(v\ne p).
\]
Explicitly, $f_vh=kf_w$ for some $k\in H_q$, so
$Ha f_vh=Ha kf_w=Ha f_w$. It is stable under multiplication by $a$
as well: $Ha\cdot a=H$, and for $v\ne q$ the braid identity gives
$Ha f_va=Hf_va f_v=Ha f_v$.
The union contains the identity, and $P(T)$ is generated by $H$ and
$a$. It is therefore all of $P(T)$, whence $[P(T):H]\leq n$ and
$|P(T)|\leq n!$. The surjection to $S_n$ is an isomorphism.
\end{proof}

\section{Transporting an exact presentation}
\begin{lemma}\label{lem:transport}
Suppose $K=\gen{s_e\mid R(s)}$. Let $r_i=v_i(s)$ generate $K$ and
choose words $w_e(r)$ with $w_e(v(s))=s_e$ in $K$. Then an exact
presentation on the marking $r_i$ is
\begin{equation}\label{eq:transport}
 \gen{r_i\mid R(w(r)),\quad r_i=v_i(w(r))\text{ for every }i}.
                                                               \end{equation}
\end{lemma}
\begin{proof}
The prescribed elements of $K$ satisfy all displayed relations, giving
a surjective map from \eqref{eq:transport} to $K$. Conversely $s_e\mapsto w_e(r)$
defines a homomorphism from $K$ to \eqref{eq:transport}, because every relator $R(w)$
is imposed. One composite fixes each $s_e$ by the word identity in
$K$; the other fixes each $r_i$ by the second group of relations.
Thus these homomorphisms are inverse.
\end{proof}


\section{Attribution and formal limits}
The existing source attributions remain in each proof and in the five
applications. The common extremal classification retains Whiston's
independent-family theorem and omission consequence separately.
Mar\'oti's primitive-group theorem and the exact finite GAP catalogue
are retained in the all-degree bound. The entire primitive-bound
program is required for its bounded exceptions; the three other
programs retain their validation and auxiliary-example roles.
All four files remain complete in the program appendix.

The existing ResidualGeneration and IrredundantNormalization Lean
files are retained unchanged. They verify the subgroup induction and
the actual normalization/independence leg, respectively. They do not
formalize laminar completion, the classification inputs, the core
attachment theorem, finite catalogue completeness, the presentations
or any full source geometry. No additional formal declaration is
claimed by this reorganization. All full costs of this common paper
and the mathematical program bundle belong to this topic.

\begin{thebibliography}{ACK05}
\bibitem[BC98]{BC98} F. Buekenhout and P. Cara,
\emph{Some properties of inductively minimal geometries},
Bull. Belg. Math. Soc. Simon Stevin \textbf{5} (1998), 213--219.
\bibitem[CC]{CC02} P. J. Cameron and P. Cara,
\emph{Independent generating sets and geometries for symmetric groups},
J. Algebra \textbf{258} (2002), 641--650. The source edition used in this
volume is the twelve-page author manuscript.
\bibitem[ACK05]{ACK05} C. Archer, P. Cara and J. Krempa,
\emph{Using the Frattini Subgroup and Independent Generating Sets to Study
RWPRI Geometries}, Beitr. Algebra Geom. \textbf{46} (2005), 169--177.
\bibitem[GL21]{GL21} N. Garcia-Colin and D. Leemans,
\emph{An infinite family of locally $X$ graphs based on incidence geometries},
Art Discrete Appl. Math. \textbf{4} (2021), \#P3.01. The exact source
edition used here is the ten-page institutional revised author manuscript,
titled \emph{An infinite family of incidence geometries whose incidence
graphs are locally $X$}.

\bibitem[FL]{FL} M. E. Fernandes and D. Leemans,
C-groups of high rank for the symmetric groups,
\emph{J. Algebra} \textbf{508} (2018), 196--218,
doi:10.1016/j.jalgebra.2018.04.031.
The edition used here is the institutional accepted manuscript,
20 authored pages plus publisher cover, not the final typeset edition.
\bibitem[M]{M} A. Mar\'oti,
On the orders of primitive groups,
\emph{J. Algebra} \textbf{258} (2002), 631--640, Theorem 1.1.
Original authored PDF retained.
\bibitem[W]{W} J. Whiston,
Maximal independent generating sets of the symmetric group,
\emph{J. Algebra} \textbf{232} (2000), 255--268,
doi:10.1006/jabr.2000.8399.
\bibitem[GAP]{GAP} The GAP Group,
\emph{GAP---Groups, Algorithms, and Programming}, version used in the
saved replay. Primitive-group catalogue and exact finite-group routines.
\end{thebibliography}
\end{document}
